\begin{circuitikz}[
        american, 
        >={Latex[length=5mm, width=2mm]},
        %>=latex,
        alt={A ladder logic diagram showing two main rungs bounded by two vertical power rails, with the left power rail highlighted in green. None of the logic elements are active or true. On the top rung, a normally open contact labeled Start PB connects to an output coil labeled LT1 on the right. On the second rung, two parallel branches each contain a normally open contact, one labeled Manual PB and the other labeled LT2. These parallel contacts join together before connecting to an output coil labeled LT2 on the right rail.}
    ]

    \def\rungs{3}   % Power rails are automatically drawn the proper height for # rungs
    \pgfmathsetmacro{\railheight}{-\rungs * 2}

    %========================================================
    % Component Grid
    % Spacing = 3x, 2y
    %========================================================
    \foreach \i in {0,...,5} {
        \foreach \j in {0,...,20} {
        
        % Calculate spatial positions
        \pgfmathsetmacro{\px}{\i * 3}
        \pgfmathsetmacro{\py}{-1 - (\j * 2)}
        
        % Name coordinate using loop indices: (C-column-row) e.g., (C-0-0), (C-5-10)
        \coordinate (C-\i-\j) at (\px, \py);
        
        % OPTIONAL: Uncomment the 2 lines below to display grid labels while building
        % \fill[red] (C-\i-\j) circle (1.5pt);
        %\node[anchor=south west, scale=0.5, gray] at (C-\i-\j) {(\i,\j)};
        }
    }

    % Put the primary logic elements on this layer
    \begin{pgfonlayer}{main}

        \draw[
            BakerBlack
        ] 
        (C-0-0) 
        to[C, color=BakerBlack, font=\small, l={Start PB}] (C-1-0) 
        -- (C-4-0) 
        to[esource, color=BakerBlack, font=\small, l={LT1}] (C-5-0);

        \draw[
            BakerBlack
        ] 
        (C-0-1)
        to[C, color=BakerBlack, font=\small, l={Manual PB}] (C-1-1)
        -- (C-4-1)
        to[esource, color=BakerBlack, font=\small, l={LT2}] (C-5-1);

        \draw[
            BakerBlack
        ]
        (C-0-2)
        to[C, color=BakerBlack, font=\small, l={LT2}] (C-1-2)
        -- (C-1-1);

    \end{pgfonlayer}

    % Use only for green highlighting of true logic elements
    \begin{pgfonlayer}{bg}
        
        % Always highlight the left rail green
        \draw[
            UpNorthGreen,
            opacity=0.6,
            line width=10pt
        ]
        (0,0)
        -- (0,\railheight);

    \end{pgfonlayer}

    % Power rails should be the only thing in the foreground
    % Automatically drawn based on the number of rungs - should need to change this
    % Rails start at y=0 and go down - note the (-) sign is in the macro above
    \begin{pgfonlayer}{ft}
        
        \draw[
            BakerBlack,
            very thick,
        ]
        (0,0)
        -- (0,\railheight);

        \draw[
            BakerBlack,
            very thick,
        ]
        (15,0)
        -- (15,\railheight);

    \end{pgfonlayer}

\end{circuitikz}